\documentclass[11pt,a4paper]{article}
\usepackage[T2A]{fontenc}
\usepackage[utf8]{inputenc}
\usepackage[russian,english]{babel}
\usepackage{amsmath,amssymb,amsthm,mathtools}
\usepackage{setspace}
\usepackage{enumitem}
\usepackage[hidelinks]{hyperref}
\onehalfspacing
\newcommand{\head}[1]{\par\medskip\noindent\textbf{\underline{#1}}\quad}
\newcommand{\sheettitle}[3]{% title, subtitle, homework-flag (empty or "Homework")
  \begin{center}
  {\huge\bfseries #1\par}\vspace{1.2em}
  {\Large #2\par}\vspace{0.8em}
  \if\relax\detokenize{#3}\relax\else{\Large #3\par}\vspace{0.8em}\fi
  {\rudate\par}
  \end{center}\vspace{0.5em}}
\newcommand{\rudate}{\foreignlanguage{russian}{\number\day~\ifcase\month\or января\or февраля\or марта\or апреля\or мая\or июня\or июля\or августа\or сентября\or октября\or ноября\or декабря\fi\ \number\year~года}}
\newcommand{\src}[1]{\hfill(#1)}
\DeclareMathOperator{\inv}{inv}
\DeclareMathOperator{\des}{des}
\DeclareMathOperator{\cyc}{cyc}
\begin{document}
\sheettitle{Combinatorics -- 8}{Permutations, binary recursions and log-concavity}{Homework}

\head{Theoretical minimum.} For $\sigma\in S_n$ one has
$\sum_\sigma q^{\inv\sigma}=\prod_{k\leqslant n}(1+\dots+q^{k-1})$ by the Lehmer code and
$\sum_\sigma x^{\cyc\sigma}=x(x+1)\cdots(x+n-1)$. The Eulerian numbers $A(n,k)$ count
permutations with $k$ descents; inserting $n$ gives
$A(n,k)=(k+1)A(n-1,k)+(n-k)A(n-1,k-1)$, and
$\sum_{m\geqslant0}(m+1)^nt^m=A_n(t)/(1-t)^{n+1}$; excedances have the same distribution
(Foata). Sign and cycle type are Abstract algebra~--~3. A complete mapping of an abelian group $G$ is a bijection $\sigma$ with
$x\mapsto x+\sigma(x)$ bijective; it forces $\sum_{g\in G}g=0$, so $\mathbb Z_n$ has one iff
$n$ is odd. Functions given by $f(2n)$, $f(2n+1)$ are handled by induction along the binary
expansion, e.g. $s_2(2n+1)=s_2(n)+1$, $v_2(n!)=n-s_2(n)$, Stern's $s(2n)=s(n)$,
$s(2n+1)=s(n)+s(n+1)$. A positive sequence with $a_k^2\geqslant a_{k-1}a_{k+1}$ (log-concave) is
unimodal, convolution preserves log-concavity, and real-rooted polynomials with non-negative
coefficients have log-concave coefficients (Newton; Polynomials~--~4). Parity:
an involution of a finite set $X$ has $|\mathrm{Fix}|\equiv|X|\pmod2$.

\medskip\noindent\rule{\linewidth}{0.4pt}

\begin{enumerate}[leftmargin=2.2em,itemsep=0.6em]
\item Prove that a log-concave sequence of positive numbers is unimodal and that the
convolution of two such sequences is log-concave. Let $B_n(t)=\sum_{m\geqslant0}(m+1)^nt^m$;
check that $B_{n+1}(t)=\bigl(tB_n(t)\bigr)'$ and prove by Rolle's theorem that $A_n(t)$ has
$n-1$ distinct negative roots. Deduce that the Eulerian numbers and the Stirling numbers
$c(n,k)$ are log-concave in $k$, and locate the maximum of $A(n,k)$.

\item Let $X$ be an arbitrary set and let $f$ be a one-to-one function mapping $X$ onto
itself. Prove that there exist mappings $g_1,g_2\colon X\to X$ such that $f=g_1\circ g_2$ and
$g_1\circ g_1=\mathrm{id}=g_2\circ g_2$, where $\mathrm{id}$ denotes the identity mapping
on $X$. \src{IMC 1997}

\item Let $N_n$ denote the number of ordered $n$-tuples of positive integers
$(a_1,a_2,\dots,a_n)$ such that $1/a_1+1/a_2+\dots+1/a_n=1$. Determine whether $N_{10}$ is
even or odd. \src{Putnam 1997 A5}

\item Define a sequence by $a_0=1$, together with the rules $a_{2n+1}=a_n$ and
$a_{2n+2}=a_n+a_{n+1}$ for each integer $n\geqslant0$. Prove that every positive rational
number appears in the set
\[
\Bigl\{\frac{a_{n-1}}{a_n}:n\geqslant1\Bigr\}=\Bigl\{\frac11,\frac12,\frac21,\frac13,\frac32,\dots\Bigr\}.
\]
(Here $a_n=s(n+1)$ is Stern's sequence.) \src{Putnam 2002 A5}

\item For every positive integer $n$ denote by $D_n$ the number of permutations
$(x_1,\dots,x_n)$ of $(1,2,\dots,n)$ such that $x_j\neq j$ for every $1\leqslant j\leqslant n$.
For $1\leqslant k\leqslant n/2$ denote by $\Delta(n,k)$ the number of permutations
$(x_1,\dots,x_n)$ of $(1,2,\dots,n)$ such that $x_i=k+i$ for every $1\leqslant i\leqslant k$ and
$x_j\neq j$ for every $1\leqslant j\leqslant n$. Prove that
\[
\Delta(n,k)=\sum_{i=0}^{k-1}\binom{k-1}{i}\frac{D_{(n+1)-(k+i)}}{n-(k+i)}.
\]
\src{IMC 2014}

\item A sequence $a_1,a_2,\dots,a_n$ is called bright if all its members are positive
integers and the numbers $a_1+1,a_2+2,\dots,a_n+n$ give all possible remainders upon
division by $n$. Let the sequence $a_1,a_2,\dots,a_{2025}$ be bright with $a_1=1$, $a_2=2$,
\dots, $a_9=9$. What is the minimal possible value of $a_1+a_2+\dots+a_{2025}$?
\src{IMC 2025, shortlist}

\item Consider a circular necklace with $2013$ beads. Each bead can be painted either white
or green. A painting of the necklace is called good if among any $21$ successive beads there
is at least one green bead. Prove that the number of good paintings of the necklace is odd.
(Two paintings that differ on some beads, but can be obtained from each other by rotating or
flipping the necklace, are counted as different paintings.) \src{IMC 2013}

\item A positive integer $n>1$ is called good if the set $\{1,2,3,\dots,3n\}$ can be divided
into ordered triples $(a_i,b_i,c_i)$, $1\leqslant i\leqslant n$, such that all $n$ quadratic
equations $a_ix^2+b_ix+c_i=0$ have a common root. Prove that there exist infinitely many
positive integers $n$ that are good and also infinitely many positive integers that are not
good. \src{IMC 2025, shortlist}

\item We say that a positive real number $d$ is good if there exists an infinite sequence
$a_1,a_2,a_3,\ldots\in(0,d)$ such that for each $n$, the points $a_1,\dots,a_n$ partition the
interval $[0,d]$ into segments of length at most $1/n$ each. Find
$\sup\{d\mid d\text{ is good}\}$. \src{IMC 2021}

\item Let $n\geqslant2$ be an integer. For a sequence $s=(s_1,\dots,s_{n-1})$ where each
$s_i=\pm1$, let $f(s)$ be the number of permutations $(a_1,\dots,a_n)$ of $\{1,2,\dots,n\}$
such that $s_i(a_{i+1}-a_i)>0$ for all $i$. For each $n$, determine the sequences $s$ for
which $f(s)$ is maximal. \src{Putnam 2025 A5}
\end{enumerate}
\end{document}
